\section{Conclusion}
\label{sec:conclusion}

We set out to select seeds under a threshold-dependent overexposure model and found that the
coverage representation behind the dominant algorithmic tool does not apply. Reverse-reachability
sampling rests on the identity $F_D(S) = c \cdot \Pr[S \cap R \neq \emptyset]$ for a random set $R$
independent of $S$, a function that is necessarily non-decreasing and submodular. On a one-hop
instance of two candidates of weight $0.4$ into one target, the model's own window distribution
gives $F_D(\{a\}) = 0.48$ against $F_D(\{a,b\}) = 0.32$, so the target is non-monotone and no such
coverage function matches it. We are explicit about the limits of that statement: it rules out the
standard coverage form and nothing more. It does not rule out signed decompositions, algorithms for
restricted instances, or Monte-Carlo estimation, and it says nothing about whether learning helps.
It also holds only for a named threshold distribution --- under uniform-threshold linear thresholds
the same instance is monotone. An earlier draft argued the stronger claim that "RR sets are empty";
that probe was mis-specified and the claim is withdrawn.

What replaces it is a measurement discipline, and an honest account of the measurements we had to
withdraw. We identified the seed \emph{fraction} $|S|/n$ as the variable that governs whether the
objective is monotone, and we reported a regime characterisation and a baseline comparison on top of
it. Both of those were produced by an activation rule that forbade an over-exposed node from turning
negative, and at Monte-Carlo budgets below the threshold the measurement itself requires; on
Congress-Twitter the corrected rule removes $60$--$63\%$ of the spread, more than every effect that
was reported. They are withdrawn and retained under an explicit marker in \S\ref{sec:experiments}.
What we can state without qualification is the constraint that caught them: in the negative-marginal
regime the label noise at $60$ trials is as large as the dispersion of the labels, the
negative-marginal share is the quantity most inflated by that noise, and no rank correlation there
should be believed below $300$ trials. We also list the eight further protocol confounds
(\S\ref{sec:limitations}) that have to be removed before a cost-versus-quality comparison between a
learned and a training-free selector is interpretable at all.

The negative result that survives is about the reference itself. The upper bound the source model
optimises is a \emph{difference} of two monotone submodular functions, and submodularity is not
preserved under differences, so even a valid upper bound would not by itself license a
sampling-based guarantee. Whether that bound is tight anywhere, and whether the marginals it must
resolve are in fact smaller than the estimation noise, are empirical questions we can no longer
answer from the measurements we have.

Two directions follow directly. The first is to re-derive the whole empirical section under a frozen
model contract --- a declared target set, seed eligibility, a budget of \emph{at most} $k$, a priced
state acquisition, and a Monte-Carlo reference with a stated tolerance --- with the eight confounds
removed. That is not a formality: two of them, the unpriced state and the weight normalisation,
interact with the model correction in the same direction. The second is to
look for a graph-level criterion --- we suspect in-/out-degree coupling --- that predicts when
degree-based heuristics fail, which would let a practitioner decide whether the overexposure
regime needs to be handled at all.
